% !TEX root = ./../../_Thesis.tex

% section's Name and Label
\section{Validation}} 
\label{sec:Validation}

To validate the visual simulation results of the refractive errors, we use a DSLR camera (Canon model EOS Rebel T3 with an 18-55mm zoom lens). The camera represents a perfect eye (\ie, without refractive aberrations). We place additional lenses in front of the camera's optical system to induce low-order aberrations (\ie, myopia, hyperopia, and astigmatism). Such lenses are placed on a support fixed to a UV filter attached to the main lens. Figure~\ref{fig:extralens} shows the camera with an additional +1.0 diopter lens attached to it. The support can hold up to three additional lenses. 
%
%results from the image filtering process, we've used additional lenses to induce low-order aberrations (\ie, myopia, hyperopia, and astigmatism) into the DSLR camera's optical system (Figure~\ref{fig:extralens}). 

For our simulations, we use a simplified eye model adjusted to the camera's settings to achieve consistent results between them.
% with the ones captured by the camera. 
More specifically, we make sure that the \emph{f-number} (\ie, the ratio of the camera lens' focal length $f$ to the diameter $D$ of its aperture):
\begin{equation}
	\centering
	\label{eq:fnumber}
	f_{number} = \frac{f}{D}
\end{equation}
\noindent
is the same for the camera and the eye model. For the experiments shown in the thesis, we fixed the focal length of the camera's main lens to 18mm (regardless of the use of additional lenses).
% to captured images (with or without extra lenses) with 18mm focal length and 
Thus, for instance, given \emph{f-number} values of $4.0$, $4.5$ and $5.0$, the corresponding camera lens aperture values are
%entrance pupil for each possible combination is
 4.5mm, 4.0mm and 3.6mm, respectively. 
Our simplified eye model (Figure~\ref{fig:syntheticeye}) has an axial diameter of 18mm. The crystalline lens causes the nodal point {\bf N} to be behind the crystalline. Thus, the eye model's effective focal length is 13.5mm: $f_{eye} = 18mm / \eta_{eye} =  18mm / 1.333 = 13.5 mm$, where $\eta_{eye}$ is the index of refraction of the eye. As a result, the eye model's pupil size (equivalent of the camera's lens aperture) needs to be rescaled to maintain the same \emph{f-number} value as the camera. Table~\ref{table:pupildiameter} shows the corresponding values of the equivalent camera apertures and pupil diameters. The simulation results shown in Chapter~\ref{chap:Results} were obtained for $f/5.0$ (third row of Table~\ref{table:pupildiameter}), although other values could have been used. 

%we must rescale the entrance pupil's diameter (Table~\ref{table:pupildiameter}) in order to maintain the same \emph{f-number} in simulations of retinal images. 

\begin{table}[h]
	\centering
	\caption{Camera apertures and pupil diameters for various f-numbers.}
	\label{table:pupildiameter}
	\begin{tabular}{rcc}
		\multicolumn{1}{l}{\bf f-number} & {\bf DSLR Camera (18mm focal length)} & {\bf Synthetic Eye (13.5mm focal length)} \\
		{\bf }               & {\bf aperture}     & {\bf pupil diameter}     \\
		{\bf $f/4.0$}         & 4.5mm                    & 3.4mm                        \\
		{\bf $f/4.5$}         & 4.0mm                    & 3.0mm                        \\
		{\bf $f/5.0$}         & 3.6mm                    & 2.7mm                        
	\end{tabular}
\end{table}

\begin{figure}[h]
	\centering
	
	\subfigure[]{
		\includegraphics[width=0.55\linewidth]{__Images/04/camera.jpg}
		\label{fig:extralens}
	}
	~
	\subfigure[]{
		\includegraphics[width=0.35\linewidth]{__Images/04/eyemodel.png}
		\label{fig:syntheticeye}
	}
	
	\caption[Optical systems used in the validation process]{Optical systems used in the validation process: (a) Canon EOS Rebel T3 with apparatus to add up to three extra lenses. Focal lens set to 18mm. (b) Simplified eye model with effective focal length of 13.5mm. {\bf N} is the nodal point.}
	\label{fig:camera}
\end{figure}